\scanpage{53}
\mspage{- 26 -}

n\textsuperscript{o} 276

on voit que le couple $(\mathrm A,\hat{\mathrm A})$ est plat.

Soient $\mathrm E$ un $\mathrm A$-module de type fini et $\mathrm F$ un sous-module de $\mathrm E$ ; identifiant $\hat{\mathrm E}$ à $\hat{\mathrm A}\handwritten{\otimes}\mathrm E$ on voit que $\hat{\mathrm F}$ est l'image de $\hat{\mathrm A}\handwritten{\otimes}\mathrm F$ dans $\hat{\mathrm E}$ ; autrement dit on a la formule
\[
\hat{\mathrm F}=\hat{\mathrm A}.\mathrm F\quad .
\]

En particulier il vient :
\[
\mathrm m^{\mathrm p}=\hat{\mathrm A}.\mathrm m^{\mathrm p}=(\hat{\mathrm A}.\mathrm m)^{\mathrm p}=(\hat{\mathrm m})^{\mathrm p}\ , 
\]
ce qui montre que la topologie de $\hat{\mathrm A}$ coincide avec la topologie $\hat{\mathrm m}$-adique Comme $\hat{\mathrm A}$ est complet et séparé pour la topologie $\hat{\mathrm m}$-adique, l'idéal $\hat{\mathrm m}$ est contenu dans le radical de $\hat{\mathrm A}$ ; comme $\mathrm G(\hat{\mathrm A})=\mathrm G(\mathrm A)$ est noethérien (parce qu'à engendrement fini sur l'anneau noethérien $\mathrm A/\mathrm m$) on voit d'autre part que $\hat{\mathrm A}$ est noethérien (\handwritten{\S}1, corollaire 2 de la proposition 7). Par suite\unsure{,} le couple $(\hat{\mathrm A},\hat{\mathrm m})$ est un anneau de Zariski, ce qui achève la démonstration.

\uline{Corollaire} - \uline{Soient} $(\mathrm A,\mathrm m)$ \uline{un anneau de Zariski et} $\mathrm E$ \uline{un A-module de type fini}. \uline{Pour tout sous-module} $\mathrm F$ \uline{de} $\mathrm E$ \uline{on a}
\[
\hat{\mathrm F}=\hat{\mathrm A}.\mathrm F\qquad\text{et }\mathrm F=\hat{\mathrm A}.\mathrm F\handwritten{\cap}\mathrm E.
\]

\uline{Si} $\mathrm F$ \uline{et} $\mathrm G$ \uline{sont deux sous-modules de} $\mathrm E$, \uline{on a}
\[
\handwritten{\widehat}{\mathrm F+\mathrm G}=\hat{\mathrm F}+\hat{\mathrm G}\ ;\quad
\handwritten{\widehat}{\mathrm F\handwritten{\cap}\mathrm G}=\hat{\mathrm F}\handwritten{\cap}\hat{\mathrm G}\ ;\quad
(\handwritten{\widehat}{\mathrm F:\mathrm G})=(\hat{\mathrm F}:\hat{\mathrm G}).
\]

La relation $\hat{\mathrm F}=\hat{\mathrm A}.\mathrm F$ a déjà été démontrée ; la relation $\mathrm F=\hat{\mathrm A}.\mathrm F\handwritten{\cap}\mathrm E$ exprime que $\mathrm F$ est fermé dans $\mathrm E$ (théorème 2) : la relation $\handwritten{\widehat}{\mathrm F+\mathrm G}=\hat{\mathrm F}+\hat{\mathrm G}$ s'écrit $\hat{\mathrm A}(\mathrm F+\mathrm G)=\unsure{.}\hat{\mathrm A}\mathrm F+\hat{\mathrm A}\mathrm G$ et est donc évidente ; pour établir la relation $\handwritten{\widehat}{\mathrm F\handwritten{\cap}\mathrm G}=\hat{\mathrm F}\handwritten{\cap}\hat{\mathrm G}$ on l'écrit $\hat{\mathrm A}.(\mathrm F\handwritten{\cap}\mathrm G)=\hat{\mathrm A}.\mathrm F\handwritten{\cap}\hat{\mathrm A}.\mathrm G$, et on applique la Prop.\unsure{,} ?? du chap.~I, \handwritten{\S} ?? sur les couples plats ; même démonstration pour la relation concernant les transporteurs.

(ici encore, le rédacteur a préféré tout bloquer dans un unique

% archive sheet number: 52
